% report/sections/03a_gaussian.tex -- T2 (Member 2): Gaussian generative classifiers (Task 2)
% Included from sections/03_models.tex (which holds \section{Models}); T3 goes in 03b_logistic.tex.
% Uses the main preamble: amsmath, booktabs and the `assumptions` environment.
% The macros below use \providecommand, so they never clash with notation.tex;
% delete them once notation.tex defines the same symbols.
\providecommand{\cFP}{c_{\mathrm{FP}}}
\providecommand{\cFN}{c_{\mathrm{FN}}}
\providecommand{\cI}{c_{\mathrm{I}}}
\providecommand{\logit}{\operatorname{logit}}
\providecommand{\tr}{\operatorname{tr}}

\subsection{Gaussian generative classifiers}\label{sec:t2}

\paragraph{Model.}
Let $x\in\mathbb{R}^d$ be the features and $y\in\{0,1\}$ the label ($y=1$ spam). The generative
model is
\begin{equation}
  P(y=1)=\pi_1,\qquad P(y=0)=\pi_0=1-\pi_1,\qquad x\mid y=k\;\sim\;\mathcal N(\mu_k,\Sigma_k).
  \label{eq:t2-model}
\end{equation}
We call it \emph{class-specific} when $\Sigma_0,\Sigma_1$ are free and \emph{shared} when
$\Sigma_0=\Sigma_1=\Sigma$. The training set is $\{(x_n,y_n)\}_{n=1}^N$ with
$N_k=\#\{n:y_n=k\}$ and $N=N_0+N_1$.

\subsubsection{Maximum-likelihood estimates}
The log-likelihood of the joint data is
\begin{equation}
  \ell=\sum_{n}\bigl[y_n\log\pi_1+(1-y_n)\log\pi_0\bigr]
      +\sum_{k=0}^{1}\sum_{n:\,y_n=k}\log\mathcal N(x_n\mid\mu_k,\Sigma_k).
  \label{eq:t2-ll}
\end{equation}

\emph{Prior.} Only the first sum involves $\pi_1$:
$\partial_{\pi_1}\bigl[N_1\log\pi_1+N_0\log(1-\pi_1)\bigr]=N_1/\pi_1-N_0/(1-\pi_1)=0$, so
\begin{equation}
  \hat\pi_1=N_1/N,\qquad \hat\pi_0=N_0/N.
\end{equation}

\emph{Means.} For class $k$, $\nabla_{\mu_k}\ell=\Sigma_k^{-1}\sum_{n:\,y_n=k}(x_n-\mu_k)=0$, hence
\begin{equation}
  \hat\mu_k=\frac1{N_k}\sum_{n:\,y_n=k}x_n .
\end{equation}

\emph{Class-specific covariances.} Put $S_k=\frac1{N_k}\sum_{n:\,y_n=k}(x_n-\hat\mu_k)(x_n-\hat\mu_k)^{\top}$
and $\Lambda_k=\Sigma_k^{-1}$. At $\mu_k=\hat\mu_k$ the class-$k$ term of \eqref{eq:t2-ll} equals
$\frac{N_k}{2}\bigl[\log|\Lambda_k|-\tr(\Lambda_kS_k)\bigr]+\mathrm{const}$. This is concave in
$\Lambda_k$, and its gradient $\frac{N_k}{2}(\Lambda_k^{-1}-S_k)$ vanishes at $\Lambda_k^{-1}=S_k$, so
\begin{equation}
  \hat\Sigma_k=S_k .
\end{equation}

\emph{Shared covariance.} With $\Sigma_0=\Sigma_1=\Sigma$ and $\Lambda=\Sigma^{-1}$ the two class terms
add to $\frac N2\bigl[\log|\Lambda|-\tr(\Lambda\,\bar S)\bigr]$ with $\bar S=(N_0S_0+N_1S_1)/N$, so the same
argument gives the pooled covariance
\begin{equation}
  \hat\Sigma=\frac{N_0\,\hat\Sigma_0+N_1\,\hat\Sigma_1}{N}
  =\frac1N\sum_{n}(x_n-\hat\mu_{y_n})(x_n-\hat\mu_{y_n})^{\top}.
  \label{eq:t2-pooled}
\end{equation}
The estimates divide by $N_k$ (or $N$), not $N_k-1$; the difference is negligible at our sample sizes.
A positive-definite $\hat\Sigma_k$ needs $N_k>d$ and non-degenerate features; see
Section~\ref{sec:t2-reg}.

\begin{assumptions}
independent samples, exactly Gaussian class-conditional densities, and $N_k>d$ with
non-degenerate features so that the estimated covariances are invertible.
\end{assumptions}

\subsubsection{Posterior: linear or quadratic in $x$}
By Bayes' rule, with $p_k(x)=\mathcal N(x\mid\mu_k,\Sigma_k)$,
\begin{equation}
  P(y=1\mid x)=\frac{\pi_1p_1(x)}{\pi_1p_1(x)+\pi_0p_0(x)}=\sigma\bigl(a(x)\bigr),
  \qquad a(x)=\log\frac{\pi_1p_1(x)}{\pi_0p_0(x)},
  \label{eq:t2-post}
\end{equation}
where $\sigma(a)=1/(1+e^{-a})$. Expanding the two Gaussian log-densities,
\begin{equation}
  a(x)=\log\frac{\pi_1}{\pi_0}-\frac12\log\frac{|\Sigma_1|}{|\Sigma_0|}
      -\frac12(x-\mu_1)^{\top}\Sigma_1^{-1}(x-\mu_1)
      +\frac12(x-\mu_0)^{\top}\Sigma_0^{-1}(x-\mu_0).
  \label{eq:t2-a}
\end{equation}

\emph{Shared covariance.} If $\Sigma_0=\Sigma_1=\Sigma$ the terms $\mp\frac12x^{\top}\Sigma^{-1}x$ cancel and
\begin{equation}
  a(x)=w^{\top}x+b,\qquad
  w=\Sigma^{-1}(\mu_1-\mu_0),\qquad
  b=\log\frac{\pi_1}{\pi_0}-\frac12\mu_1^{\top}\Sigma^{-1}\mu_1+\frac12\mu_0^{\top}\Sigma^{-1}\mu_0 ,
  \label{eq:t2-linear}
\end{equation}
so $a$ is \emph{linear} in $x$.

\emph{Class-specific covariances.} The quadratic terms no longer cancel:
\begin{equation}
  a(x)=\tfrac12x^{\top}\bigl(\Sigma_0^{-1}-\Sigma_1^{-1}\bigr)x
      +\bigl(\Sigma_1^{-1}\mu_1-\Sigma_0^{-1}\mu_0\bigr)^{\top}x+c,
  \label{eq:t2-quad}
\end{equation}
with $c=\log\frac{\pi_1}{\pi_0}-\frac12\log\frac{|\Sigma_1|}{|\Sigma_0|}
-\frac12\mu_1^{\top}\Sigma_1^{-1}\mu_1+\frac12\mu_0^{\top}\Sigma_0^{-1}\mu_0$, so $a$ is \emph{quadratic}
in $x$ whenever $\Sigma_0\neq\Sigma_1$.

\begin{assumptions}
the model \eqref{eq:t2-model} holds; the fitted posterior is \eqref{eq:t2-post}
with the plug-in estimates.
\end{assumptions}

\subsubsection{Decision boundary for any probability threshold}
A rule ``send to spam iff $P(y=1\mid x)\ge t$'' is equivalent to $a(x)\ge\logit(t)=\log\frac{t}{1-t}$,
because $\sigma$ is strictly increasing. The boundary is therefore the level set
\begin{equation}
  \{x:\;a(x)=\log\tfrac{t}{1-t}\}.
  \label{eq:t2-boundary}
\end{equation}
For the shared covariance this is the hyperplane $w^{\top}x=\log\frac{t}{1-t}-b$. Its normal $w$ does not
depend on $t$: changing $t$ only translates the hyperplane, and $t=\frac12$ gives $w^{\top}x+b=0$. For
class-specific covariances it is a quadric surface (an ellipsoid, paraboloid, hyperboloid or a pair of
hyperplanes, depending on the signs of the eigenvalues of $\Sigma_0^{-1}-\Sigma_1^{-1}$).

With the loss matrix of Task~1 the relevant thresholds are $\theta_1=1-\cI/\cFP$ and $\theta_0=\cI/\cFN$,
and the two-action threshold is $t=\cFP/(\cFP+\cFN)$. In logit form
\begin{equation}
  \logit(t)=\log\frac{\cFP}{\cFN},\qquad
  \logit(\theta_1)=\log\frac{\cFP-\cI}{\cI},\qquad
  \logit(\theta_0)=\log\frac{\cI}{\cFN-\cI}.
  \label{eq:t2-logit-thresholds}
\end{equation}
So for the shared covariance the inspect region is the slab between two parallel hyperplanes
$a(x)=\logit(\theta_0)$ and $a(x)=\logit(\theta_1)$, and moving the cost ratio $\cFP/\cFN$ only shifts the
boundary along $w$.

\begin{assumptions}
the posterior $\sigma(a(x))$ is correctly calibrated, so that thresholding it at the
Task~1 value is the Bayes rule; otherwise \eqref{eq:t2-boundary} is only a heuristic.
\end{assumptions}

\subsubsection{Covariance regularisation}\label{sec:t2-reg}
\paragraph{Why Spambase needs it.}
Spambase has $d=57$ features and, with a 60/20/20 split, about $2761$ training emails. A class-specific
model has $2\cdot\frac{d(d+1)}2=3306$ covariance parameters, more than there are training emails, and
the shared model has $1653$. Many word-frequency features are zero for most emails, so within a class
they are nearly or exactly constant. Then $\hat\Sigma_k$ has (near-)zero eigenvalues,
$\log|\hat\Sigma_k|\to-\infty$ and $\hat\Sigma_k^{-1}$ is huge, so the quadratic form in \eqref{eq:t2-a}
is dominated by a few directions and the posteriors saturate at $0$ or $1$. This is harmful for
cost-sensitive decisions, which depend on where $p$ sits relative to $\theta_0,\theta_1$.
\emph{[Insert the measured smallest eigenvalue and condition number of $\hat\Sigma_k$ on the training
split.]}

\paragraph{Rule.}
We shrink each estimated covariance towards a scaled identity:
\begin{equation}
  \tilde\Sigma(\lambda)=(1-\lambda)\,\hat\Sigma+\lambda\,\bar v\,I,\qquad
  \bar v=\frac{\tr(\hat\Sigma)}{d},\qquad\lambda\in[0,1],
  \label{eq:t2-shrink}
\end{equation}
applied to $\hat\Sigma$ in the shared model and to each $\hat\Sigma_k$ in the class-specific model. If
$\hat\Sigma$ has eigenvalues $\nu_i\ge0$ then $\tilde\Sigma$ has eigenvalues $(1-\lambda)\nu_i+\lambda\bar v$,
so for any $\lambda>0$ and $\hat\Sigma\ne0$ they are all at least $\lambda\bar v>0$ and the estimate is
positive definite. The trace (average variance) is preserved, $\lambda=0$ recovers the MLE, and
$\lambda=1$ gives a spherical model. The strength $\lambda$ is chosen on the validation set from a grid,
by minimum validation log-loss; the test set is never used. Features are standardised on the
training data first, so the identity is a sensible target scale.

\paragraph{Numerics.}
With $\tilde\Sigma_k=L_kL_k^{\top}$ (Cholesky), $\log|\tilde\Sigma_k|=2\sum_i\log(L_k)_{ii}$ and
$(x-\mu_k)^{\top}\tilde\Sigma_k^{-1}(x-\mu_k)=\|L_k^{-1}(x-\mu_k)\|^2$ by a triangular solve, so no explicit
inverse is formed. The posterior is $\sigma(a(x))$ with $a(x)=\log\pi_1p_1-\log\pi_0p_0$ computed from
log-densities.

\begin{assumptions}
shrinkage towards a scaled identity is an adequate bias, features are on comparable
scales, and validation log-loss tracks test performance.
\end{assumptions}

\subsubsection{Link to logistic regression}\label{sec:t2-lr}
For the shared covariance, \eqref{eq:t2-linear} shows that the Gaussian model has exactly the posterior
form of logistic regression, $P(y=1\mid x)=\sigma(w^{\top}x+b)$. The Gaussian parameters that play the role
of the weights and bias are
\[
  w=\Sigma^{-1}(\mu_1-\mu_0),\qquad
  b=\log\frac{\pi_1}{\pi_0}-\tfrac12(\mu_1+\mu_0)^{\top}\Sigma^{-1}(\mu_1-\mu_0).
\]
(The two forms of $b$ agree because $-\frac12\mu_1^{\top}\Sigma^{-1}\mu_1+\frac12\mu_0^{\top}\Sigma^{-1}\mu_0
=-\frac12(\mu_1+\mu_0)^{\top}\Sigma^{-1}(\mu_1-\mu_0)$ for symmetric $\Sigma^{-1}$.) The class prior enters only through
$b$, so a change of prior from $\pi_1$ to $\pi_1'$ shifts the intercept by
$\log\frac{\pi_1'/\pi_0'}{\pi_1/\pi_0}$; this is the prior-shift correction of Task~4.
For class-specific covariances, \eqref{eq:t2-quad} is a logistic regression on the quadratic feature map
$\bigl(x,\ \mathrm{vech}(xx^{\top})\bigr)$.

The two methods differ in how the parameters are fitted and how many are estimated. The Gaussian model
maximises the \emph{joint} likelihood \eqref{eq:t2-ll}, whereas logistic regression maximises the
\emph{conditional} likelihood of $y$ given $x$ and leaves $(w,b)$ free (Task~2, T3). Counting parameters for $d=57$:

\begin{center}
\begin{tabular}{lcc}
\toprule
Model & Free parameters & $d=57$\\
\midrule
Gaussian, shared covariance & $2d+\tfrac{d(d+1)}2+1$ & $1768$\\
Logistic regression (linear) & $d+1$ & $58$\\
Gaussian, class-specific covariances & $2d+d(d+1)+1$ & $3421$\\
Logistic regression on quadratic features & $d+\tfrac{d(d+1)}2+1$ & $1711$\\
\bottomrule
\end{tabular}
\end{center}

If the class-conditionals really are Gaussian with a common covariance, the generative fit uses the
structure and is the more efficient estimator. If they are not (for example the heavy-tailed, mostly zero
Spambase features), logistic regression still targets $P(y\mid x)$ directly and is usually the more
robust choice. We therefore expect different calibration for the two models, which is what the
cost-sensitive decisions depend on.

\begin{assumptions}
for the equivalence, a common covariance $\Sigma$; the efficiency comparison also
assumes the Gaussian model is correctly specified.
\end{assumptions}
